% Figure: event-prefix balances and their sum; separate against pooled peak stock. Data from model/grid.ts.
\begin{figure}[t]
\centering
\begin{tikzpicture}
\begin{axis}[
  width=0.96\textwidth, height=6.6cm,
  xlabel={round}, ylabel={balance (sat)},
  xmin=0, xmax=\walkT, ymin=0,
  scaled y ticks=false, yticklabel style={/pgf/number format/fixed, /pgf/number format/1000 sep={\,}},
  grid=major, grid style={black!12},
  tick label style={font=\footnotesize}, label style={font=\small},
  legend style={font=\footnotesize, draw=none, fill=none, at={(0.5,-0.2)}, anchor=north, legend columns=2, /tikz/every even column/.append style={column sep=0.6cm}},
  legend cell align=left,
]
\addplot[black!35, thin, const plot] table[x=t, y=a1] {\walkpaths}; \addlegendentry{one agent's balance (four shown)}
\addplot[black!35, thin, const plot, forget plot] table[x=t, y=a2] {\walkpaths};
\addplot[black!35, thin, const plot, forget plot] table[x=t, y=a3] {\walkpaths};
\addplot[black!35, thin, const plot, forget plot] table[x=t, y=a4] {\walkpaths};
\addplot[black, thick, const plot] table[x=t, y=sum] {\walkpaths}; \addlegendentry{the four together}
\addplot[black!60, dashed, thick, domain=0:\walkT, samples=2] {\walkpeaksum}; \addlegendentry{sum of the four peaks: per-agent stock}
\addplot[black, dotted, thick, domain=0:\walkT, samples=2] {\walksumpeak}; \addlegendentry{peak of the sum: pooled stock}
\end{axis}
\end{tikzpicture}
\caption{Four agents with zero offered drift, drawn with the model's generator over \walkT{} rounds. Each receipt precedes that agent's spend; an unaffordable spend is rejected in full. Fractional positions within a round expose this event order. The paths show the resulting finite jump processes, with every state-changing event prefix retained. The dashed line is the sum of individual event-prefix peaks; the dotted line is the peak of aggregate holdings. These are the separate and pooled oracle stocks for this trace with zero safety margin. Their ratio measures this trace's pooling gain; the diffusion comparison in Remark~\ref{rem:sqrtN} uses additional assumptions. Generated by \texttt{model/grid.ts}.}
\label{fig:walks}
\end{figure}
